\chapter{ОПИСАНИЕ ПРЕДЛАГАЕМОГО МЕТОДА}

В этой главе описана архитектура адаптивного токенизатора ДНК-последовательностей. Метод основан на рекурсивном гейтированном слиянии с дифференцируемым динамическим программированием и обучается совместно с backbone-моделью на задаче предсказания следующего нуклеотида (NTP) на референсном геноме человека.

\section{Общая архитектура}

Архитектура состоит из четырёх этапов:
\begin{enumerate}
  \item инициализация представлений на уровне нуклеотидов;
  \item рекурсивное слияние соседних элементов с гейтированной оценкой;
  \item кодирование результата backbone-моделью на укороченной последовательности;
  \item декодирование обратно на уровень нуклеотидов для вычисления NTP loss.
\end{enumerate}

Идея здесь простая: начинаем с отдельных нуклеотидов, рекурсивно объединяем соседей в составные токены, а основную вычислительную работу выполняем уже на более короткой последовательности. Какие именно пары сливать, решает обучаемая функция, оптимизируемая вместе с моделью.

От фиксированного patching (MegaByte~\cite{yuMegaBytePredictingMillionbyte2023}, PatchDNA~\cite{patchDNA2026}) метод отличается адаптивными границами. От одноуровневой мягкой токенизации (GBST в Charformer~\cite{tayCharformerFastCharacter2022}) - рекурсивной схемой с ограничением непересечения через дифференцируемое DP. От entropy-based patching (BLT~\cite{bltByteLatentTransformer2024}) - использованием биологических приоров. От MergeDNA~\cite{mergeDNAContextAware2026} - рекурсией и дифференцируемым DP.


\section{Инициализация представлений}

\subsection{Входная последовательность и embedding}

Входная ДНК-последовательность длины $L$:
\begin{equation}
S = (n_1, n_2, \dots, n_L), \qquad n_i \in \{A, T, C, G\}.
\label{eq:input_seq}
\end{equation}

Словарь включает канонические нуклеотиды, символ $N$ и специальные токены (CLS, SEP, PAD, MASK, UNK) - всего $|V| = 10$. Каждый входной символ отображается в вектор размерности $d$ через обучаемую матрицу $E \in \mathbb{R}^{|V| \times d}$:
\begin{equation}
e_v = E[v] \in \mathbb{R}^d.
\label{eq:embedding}
\end{equation}

Словарь крошечный по меркам NLP (там порядка $10^4$--$10^5$ токенов), и вся выразительная способность представления переносится на последующие уровни слияния.

Инициализацию embedding matrix можно сделать биологически мотивированной: пурины $(A, G)$ и пиримидины $(T, C)$ образуют две группы, а комплементарные пары $(A \leftrightarrow T, G \leftrightarrow C)$ связаны Watson--Crick спариванием. На практике каждому нуклеотиду назначается вектор группы, затем добавляется уникальное смещение и небольшой случайный шум для нарушения симметрии.

\subsection{Позиционное кодирование}

Используется Rotary Positional Embedding (RoPE)~\cite{suRoFormerEnhancedTransformer2021}~- скалярное произведение двух закодированных позиций зависит только от относительного расстояния. Это полезно для рекурсивного слияния, где абсолютные позиции теряют смысл, а относительное положение соседей остаётся важным.

\subsection{Контекстное обогащение}

Начальные embedding'и содержат информацию только об одном нуклеотиде. Для расширения контекста перед merge pipeline применяется каузальная одномерная свёртка с kernel size~5:
\begin{equation}
\hat{x}_i^{(0)} = x_i^{(0)} + \operatorname{Conv1D}_{k=5}(\mathbf{X}^{(0)})_i.
\label{eq:contextualizer}
\end{equation}

Каузальность (padding слева) нужна для совместимости с NTP. При этом параметров добавляется $5d^2$ - при $d = 256$ около 320\,000, что пренебрежимо по сравнению с backbone.

\subsection{Обработка reverse complement}

ДНК - двуцепочечная молекула, и многие функциональные элементы работают на обеих цепях. Поэтому используются две стратегии: data augmentation (с вероятностью $0{,}5$ заменяем последовательность на reverse complement) и опциональный регуляризатор $\mathcal{L}_{\mathrm{RC}}$, штрафующий расхождение оценок слияния между прямой и комплементарной цепью (подробнее в~\S\ref{sec:loss}).


\section{Гейтированная оценка границ}

\subsection{Функция оценки слияния}

На каждом уровне рекурсии $l$ для каждой соседней пары $(x_i^{(l)}, x_{i+1}^{(l)})$ вычисляется оценка слияния из двух компонент.

Оценка сходства - scaled dot-product с обучаемой проекцией:
\begin{equation}
s_{\mathrm{sim},i}^{(l)} = \frac{(x_i^{(l)})^\top W_k^{(l)}\, x_{i+1}^{(l)}}{\sqrt{d}},
\label{eq:sim_score}
\end{equation}
где $W_k^{(l)} \in \mathbb{R}^{d \times d}$. Проекция делает оценку асимметричной, что не недостаток: ДНК-последовательность направлена ($5' \to 3'$), и порядок нуклеотидов биологически значим~\cite{rosenbergPatternsTransitionalMutation2003}.

Гейт - контекстно-зависимая склонность к слиянию:
\begin{equation}
g_i^{(l)} = \sigma\!\left((w_g^{(l)})^\top [x_i^{(l)}; x_{i+1}^{(l)}] + b_g^{(l)}\right),
\label{eq:gate}
\end{equation}
где $\sigma$ - сигмоида. Гейт может подавлять слияние на границах кодонов или в регуляторных областях даже при высоком сходстве соседей. Смещение $b_g^{(l)}$ инициализируется $-0{,}5$, чтобы на старте модель предпочитала не сливать (предотвращает чрезмерное сжатие на ранних итерациях).

Итого: $s_i^{(l)} = g_i^{(l)} \cdot s_{\mathrm{sim},i}^{(l)}$.

\subsection{Температурная модуляция}

Для управления жёсткостью решений вводится температура $\tau^{(l)} > 0$:
\begin{equation}
s_i^{(l)} = g_i^{(l)} \cdot \frac{s_{\mathrm{sim},i}^{(l)}}{\tau^{(l)}}.
\label{eq:temp_modulation}
\end{equation}

Температура снижается по cosine schedule в процессе обучения: от $\tau_{\max} = 1{,}0$ до $\tau_{\min} = 0{,}1$. Высокая температура означает мягкие решения, низкая - жёсткие. Это обеспечивает плавный переход от мягких решений при обучении к дискретным при инференсе.

\subsection{Интеграция биологических приоров}

Биологические приоры вводятся как дополнительные входы в гейт, а не как отдельная цель обучения.

Для оценок консервативности (phastCons~\cite{siepelEvolutionarilyConservedElements2005}, $c_i \in [0,1]$) формируется вектор из пяти признаков: значения $c_i$ и $c_{i+1}$, их разность, минимум и среднее. Разрыв $|c_i - c_{i+1}|$ может сигнализировать о функциональной границе.

Для мотивов JASPAR~\cite{rauluseviciuteJASPAR2024} аналогично: индикаторы начала и конца мотива, число мотивов на позиции, качество совпадения.

Расширенный гейт объединяет обучаемую и биологическую компоненты:
\begin{equation}
g_i^{(l)} = \sigma\!\left(
(w_g^{(l)})^\top [x_i^{(l)}; x_{i+1}^{(l)}]
+ (w_c^{(l)})^\top \phi_i
+ (w_m^{(l)})^\top \psi_i
+ b_g^{(l)}
\right),
\label{eq:extended_gate}
\end{equation}
где $\phi_i$ и $\psi_i$ - векторы признаков консервативности и мотивов (по 5 компонент каждый), $w_c^{(l)}$ и $w_m^{(l)}$ - обучаемые веса. Если аннотации отсутствуют, модель использует только первый член.


\section{Механизм рекурсивного слияния}

\subsection{Задача выбора пар}

На каждом уровне нужно выбрать подмножество соседних пар для слияния, причём выбранные пары не должны пересекаться:
\begin{equation}
M^{(l)} = \arg\max_{M} \sum_{i \in M} s_i^{(l)},
\qquad |M| = \lfloor r^{(l)} N^{(l)} \rfloor,
\qquad \forall\, i, j \in M: |i - j| > 1,
\label{eq:discrete_selection}
\end{equation}
где $r^{(l)}$ - целевая доля слияний на уровне $l$. Задача решается динамическим программированием на одномерной последовательности, однако жёсткий выбор делает этот шаг недифференцируемым.

\subsection{Дифференцируемый выбор через dDP}
\label{sec:ddp}

При обучении жёсткий argmax заменяется на мягкий через log-semiring relaxation~\cite{menschDifferentiableDynamicProgramming2018}. Вместо операции $\max$ используется $\mathrm{softmax}_\beta$:
\begin{equation}
\mathrm{softmax}_\beta(a, b) = \frac{1}{\beta} \log(\exp(\beta a) + \exp(\beta b)),
\label{eq:softmax_beta}
\end{equation}
которая при $\beta \to \infty$ стремится к $\max$.

DP-рекуррента принимает вид:
\begin{equation}
D[i, k] = \mathrm{softmax}_\beta\!\left(
D[i-1, k],\;
D[i-2, k-1] + s_i^{(l)}
\right),
\label{eq:dp_recurrence}
\end{equation}
где $D[i, k]$ - мягкий оптимальный score при $i$ рассмотренных позициях и $k$ выбранных парах. Мягкие merge-веса $\alpha_i^{(l)} \in [0,1]$ извлекаются как производная $\partial D[N, K] / \partial s_i$, что интерпретируется как маргинальная вероятность включения позиции $i$ в оптимальный набор. Сложность - $O(N \cdot K)$ на уровень.

В качестве альтернативы рассматривается Straight-Through Estimator (STE)~\cite{bengioEstimatingPropagatingGradients2013}: forward pass - жёсткий DP, backward - градиент через soft scores. Подход проще в реализации, но даёт смещённые градиенты.

\subsection{Формирование объединённых представлений}

Для каждой выбранной пары объединённое представление вычисляется MLP:
\begin{equation}
\tilde{x}_{i}^{(l+1)} = \operatorname{MLP}_{\mathrm{merge}}^{(l)}\!\left([x_i^{(l)};\; x_{i+1}^{(l)};\; x_i^{(l)} \odot x_{i+1}^{(l)}]\right),
\label{eq:merge_mlp}
\end{equation}
где $\odot$ - поэлементное произведение. Конкатенация сохраняет информацию, произведение кодирует попарное взаимодействие. MLP: вход $3d$, скрытый $2d$, выход $d$, активация GELU~\cite{hendrycksGaussianErrorLinear2016}.

При обучении переход к следующему уровню - мягкая интерполяция:
\begin{equation}
x_i^{(l+1)} = \alpha_i^{(l)} \cdot \tilde{x}_{i}^{(l+1)} + (1 - \alpha_i^{(l)}) \cdot \operatorname{proj}(x_i^{(l)}),
\label{eq:soft_assignment}
\end{equation}
где $\operatorname{proj}$ - линейная проекция с LayerNorm. При инференсе $\alpha_i \in \{0, 1\}$ - жёсткий выбор.

\subsection{Динамика сжатия}

После каждого уровня длина сокращается. При долях слияний $r^{(1)} = 0{,}3$, $r^{(2)} = 0{,}3$, $r^{(3)} = 0{,}25$ итоговое сжатие составляет $1 - 0{,}7 \times 0{,}7 \times 0{,}75 \approx 63\%$: из 2048 нуклеотидов получается около 756 токенов. Чтобы спаны не становились слишком длинными, вводится ограничение $\ell_j \leq \ell_{\max} = 8$: пары, слияние которых нарушило бы ограничение, исключаются из кандидатов.

\subsection{Training--inference gap}

Мягкий selection при обучении и жёсткий при инференсе создают расхождение. Чтобы его контролировать, отслеживается доля <<неуверенных>> позиций ($\alpha_i \in (0{,}2;\, 0{,}8)$), несовпадение soft и hard решений, а также разность val loss между режимами. Если разность превышает $0{,}05$, температурный annealing ускоряется; если меньше $0{,}01$, наоборот, замедляется.


\section{Позиционное кодирование после слияния и декодирование}

\subsection{Span positional encoding}

После слияния каждый токен $j$ покрывает непрерывный диапазон $[a_j, b_j]$ в исходной последовательности. RoPE частично размывается после нелинейных MLP, поэтому перед backbone повторно внедряется позиционная информация.

Для каждого токена формируется вектор признаков из координат начала и конца спана, его длины, центра и нормализованных значений, который затем пропускается через MLP. Дополнительно кодируется глубина рекурсии через one-hot.

\subsection{Backbone}

Укороченная последовательность ($N \ll L$ токенов) обрабатывается каузальным Transformer: 6 слоёв, 8 голов, $d = 256$, FFN $4d = 1024$, pre-LN, RoPE на уровне token-позиций.

\subsection{Cross-attention каузальный декодер}
\label{sec:decoder}

Целевая функция определена на уровне нуклеотидов, поэтому нужно восстановить представления длины $L$. Проблема в каузальности: токен $j$ со спаном $[a_j, b_j]$ содержит информацию обо всех нуклеотидах внутри, включая <<будущие>>. Следовательно, декодер должен это учитывать явно.

Для каждой позиции $t$ формируется query из embedding'а предыдущего нуклеотида. Каузальная маска допускает attention только к токенам с $b_j < t$ - полностью завершившимся до позиции $t$. Residual skip connection от shifted embeddings даёт декодеру прямой доступ к nucleotide-level информации:
\begin{equation}
r_t = \operatorname{CrossAttn}(q_t, \hat{Z}, \operatorname{Mask}(t, \cdot)) + W_{\mathrm{skip}} \cdot E[n_{t-1}].
\label{eq:residual_decoder}
\end{equation}

На выходе - логиты по четырём нуклеотидам через линейный слой + softmax.

Для длинных последовательностей ($L > 4096$) предусмотрена альтернатива: intra-span causal decoder, где нуклеотиды внутри каждого спана предсказываются авторегрессионно через лёгкий GRU.


\section{Функция потерь}
\label{sec:loss}

\subsection{Основной loss}

Next-nucleotide prediction:
\begin{equation}
\mathcal{L}_{\mathrm{NTP}} = -\frac{1}{|T_{\mathrm{valid}}|}\sum_{t \in T_{\mathrm{valid}}} \log P(n_t \mid n_{<t}),
\label{eq:ntp_loss}
\end{equation}
где $T_{\mathrm{valid}}$ исключает позиции с $N$ и PAD. Границы токенов обучаются через то, как они помогают backbone и декодеру решать NTP.

\subsection{Регуляризаторы}

К основному loss добавляются четыре регуляризатора.

Compression loss штрафует отклонение фактического коэффициента сжатия от целевого на каждом уровне:
\begin{equation}
\mathcal{L}_{\mathrm{compress}} = \sum_{l=1}^{L_{\max}} \left( \frac{|\mathcal{M}^{(l)}|}{N^{(l)}} - r_{\mathrm{target}}^{(l)} \right)^2.
\label{eq:compression_loss}
\end{equation}

Entropy regularization поощряет exploration решений на ранних стадиях. Этот член используется только в warmup-фазе, затем обнуляется.

RC-consistency regularization штрафует расхождение merge-скоров между прямой и комплементарной цепями:
\begin{equation}
\mathcal{L}_{\mathrm{RC}} = \frac{1}{|\mathcal{P}|}\sum_{(i, \mathrm{RC}(i)) \in \mathcal{P}} \left( s_i^{(l)} - s_{\mathrm{RC}(i)}^{(l)} \right)^2.
\label{eq:rc_loss}
\end{equation}

Span diversity loss штрафует низкую дисперсию длин токенов - если все спаны одинаковой длины, токенизация по сути не адаптивна.

\subsection{Итого}

\begin{equation}
\mathcal{L} = \mathcal{L}_{\mathrm{NTP}} + \lambda_{\mathrm{comp}} \mathcal{L}_{\mathrm{compress}} + \lambda_{\mathrm{ent}} \mathcal{L}_{\mathrm{ent}} + \lambda_{\mathrm{RC}} \mathcal{L}_{\mathrm{RC}} + \lambda_{\mathrm{div}} \mathcal{L}_{\mathrm{div}}.
\label{eq:total_loss}
\end{equation}

Коэффициенты $\lambda$ меняются по расписанию. В warmup (первые 10\% шагов): $\lambda_{\mathrm{comp}} = 1{,}0$, $\lambda_{\mathrm{ent}} = 0{,}1$, $\lambda_{\mathrm{div}} = 0{,}01$, $\lambda_{\mathrm{RC}} = 0{,}1$. В основной фазе $\lambda_{\mathrm{ent}}$ обнуляется, $\lambda_{\mathrm{comp}}$ снижается до $0{,}1$ по cosine schedule. В фазе hardening всё фиксируется.


\section{Стратегия обучения}

\subsection{Optimizer}

AdamW~\cite{loshchilovDecoupledWeightDecay2019} ($\beta_1 = 0{,}9$, $\beta_2 = 0{,}98$, weight decay $0{,}01$), пиковый learning rate $3 \times 10^{-4}$, linear warmup 2000 steps, cosine decay до $3 \times 10^{-5}$. Gradient clipping: max norm $1{,}0$. bf16 mixed precision.

\subsection{Phased training}

Обучение состоит из трёх фаз, и переход между ними задаётся расписанием.

Phase~1 (backbone warmup, steps 0--5\,000): merge pipeline заморожен, слияний нет, backbone учится как character-level модель. Цель - получить осмысленные embedding'и до того, как merge pipeline начнёт работать.

Phase~2 (joint training, steps 5\,000--90\,000): merge pipeline разморожен. Температурный annealing и рост $\beta$ активированы. Коэффициенты $\lambda$ по расписанию.

Phase~3 (hardening, steps 90\,000--100\,000): финальный annealing до почти дискретных решений.

\subsection{Data pipeline}

Референсный геном человека CHM13~\cite{nurgalievT2TCHm13CompleteHuman2022}; разбиение хромосом задаётся по протоколу HyenaDNA~\cite{nguenHyenaDNALongRangeGenomic2024}: chr8 - test, chr9 - validation. Далее берутся случайные фрагменты длины $L$; регионы с $>50\%$ $N$ пропускаются. RC augmentation применяется с вероятностью $0{,}5$. Conservation scores получаются через BigWig (phastCons~\cite{siepelEvolutionarilyConservedElements2005}), motif annotations - через BigBed (JASPAR~\cite{rauluseviciuteJASPAR2024}). Batch size 64, gradient accumulation $\times 2$ (effective batch 128).


\section{План экспериментальной оценки}

\subsection{Baseline модели}

Для сравнения: Char-Transformer (character-level), BPE-Transformer (словарь 4096), Fixed-kmer-Transformer ($k = 6$), Char-Mamba (SSM), PatchDNA~\cite{patchDNA2026}, Charformer~\cite{tayCharformerFastCharacter2022}. Все обучаются на тех же данных с сопоставимым числом параметров ($\pm 10\%$).

\subsection{Ablation study}

Отключаем по одному компоненту: bio priors, температурный annealing, span PE, произведение Адамара в MLP merge, dDP (замена на STE), cross-attention декодер (замена на MLP), контекстуализатор, phased training.

\subsection{Метрики}

Предобучение: NTP perplexity (soft и hard modes), bits per nucleotide, effective compression ratio, merge decision entropy, training--inference gap.

Downstream: DART-Eval~\cite{patel2025dartevalcomprehensivednalanguage} (zero-shot, probing, fine-tuning). Fine-tuning: backbone frozen или LoRA rank~8, linear probe, 10 epochs, lr $10^{-4}$. Merge pipeline при fine-tuning не дообучается. Дополнительно: variant effect prediction (ClinVar~\cite{landrumClinVar2018}), promoter detection (EPDnew~\cite{dreosEPDnew2015}), splice site prediction (GENCODE~\cite{frankishGENCODE2019}).

\subsection{Анализ выученной токенизации}

Для проверки биологической гипотезы: enrichment analysis (пересечение границ с аннотациями GENCODE~\cite{frankishGENCODE2019} и ENCODE cCREs~\cite{encodeExpandedRegistry2024}), biological coherence score (совпадение с exon-intron junctions и границами мотивов TF), корреляция длины спана с conservation score, визуализация в genome browser.


\section{Вычислительная сложность}

Основные компоненты: embedding - $O(Ld)$; merge scoring - $O(N^{(l)} d^2)$ на уровень; dDP - $O(N^{(l)} K^{(l)})$ на уровень; merge MLP - $O(N^{(l)} d^2)$ на уровень; backbone Transformer - $O(N^2 d)$; cross-attention decoder - $O(LNd)$.

Суммарно около 19{,}5M параметров (backbone ${\sim}15$M, merge pipeline ${\sim}3$M, decoder ${\sim}1{,}5$M). При этом самый дорогой компонент (backbone) работает на сжатой последовательности $N \ll L$.


\section{Выводы по главе~2}

Предложена архитектура адаптивного токенизатора, отличающаяся от существующих работ комбинацией рекурсивного слияния, дифференцируемого DP через log-semiring и интеграции биоприоров в гейтирование. Температурный annealing обеспечивает переход от мягких решений при обучении к жёстким при инференсе. Каузальный cross-attention декодер восстанавливает нуклеотидное разрешение для вычисления NTP loss.

Описаны составная функция потерь с четырьмя регуляризаторами, трёхфазная стратегия обучения и план экспериментов на DART-Eval с ablation study.
